% ============================================================
%  Outage-Aware Campus Microgrid Control with Deep RL
%  2-page report  —  IEEE two-column format
% ============================================================
\documentclass[10pt,twocolumn]{article}

% ---- packages -----------------------------------------------
\usepackage[margin=1.8cm,top=2cm,bottom=2cm]{geometry}
\usepackage{times}
\usepackage{amsmath,amssymb}
\usepackage{booktabs}
\usepackage{graphicx}
\usepackage{microtype}
\usepackage{hyperref}
\usepackage{xcolor}
\usepackage{caption}
\captionsetup{font=small,skip=4pt}

% ---- title block --------------------------------------------
\title{\textbf{Outage-Aware Campus Microgrid Control\\
       with Deep Reinforcement Learning}}
\author{[Name 1, Name 2, \ldots]\quad
        \textit{AI, Assignment 2 --- Real-Life Use Cases of RL}}
\date{}

% =============================================================
\begin{document}
\maketitle
\thispagestyle{empty}

% ---- ABSTRACT -----------------------------------------------
\begin{abstract}
We train a family of DQN agents---vanilla, Double, Dueling,
Prioritised-Experience-Replay (PER), and Lagrangian-constrained
variants---to operate a simulated campus solar-plus-battery microgrid.
The best unconstrained agent (Dueling Double DQN + PER) cuts
\textbf{operating cost (energy purchases plus comfort and wear
penalties) by about 17\%} relative to the hand-written rule
baseline and closes 80\,\% of the gap to a perfect-foresight oracle,
while a Lagrangian multiplier scheme is investigated for critical-load
protection. All components are implemented from scratch in PyTorch.
\end{abstract}

% =============================================================
\section{Problem and RL Formulation}

Each \textit{episode} is one simulated day (24 hourly steps).
The \textbf{state} is an 11-dimensional vector: battery state of charge
(\%), hour encoded as $(\sin,\cos)$, current solar generation, a 3-hour
noisy solar forecast, electricity price, total load, a binary
outage-active flag, and a noisy outage-warning signal.
The \textbf{action space} has six discrete actions:
\{charge, discharge, idle\}~$\times$~\{keep, shed non-critical load\}.

The \textbf{reward} is
\[
r_t = -\bigl(c_{\text{grid}} + c_{\text{shed}} + c_{\text{cut}}
             + c_{\text{wear}} + \lambda\,u_{\text{crit}}\bigr)\big/1000,
\]
where $c_{\text{grid}}$ is grid energy purchase cost, $c_{\text{shed}}$
is voluntary shedding, $c_{\text{cut}}$ is forced-cut penalty,
$c_{\text{wear}}$ is battery-cycle wear, and $u_{\text{crit}}$ is
unserved critical energy.  For constrained agents $\lambda$ is updated
by dual ascent so that average $u_{\text{crit}}$ stays near a per-day
budget.

\textbf{Data.} Real hourly solar irradiance from NASA POWER (2018--2023).
Split: train 2018--2021, validation 2022, held-out test 2023.
Battery: 400\,kWh / 100\,kW; 30\,\% of days have outages (2--4\,h).
All load, tariff, and outage parameters are simulated assumptions.

% =============================================================
\section{Methods}

We implement four incremental DQN improvements applied cumulatively:

\begin{enumerate}\setlength\itemsep{1pt}
  \item \textbf{Vanilla DQN} --- fixed target network, replay buffer.
  \item \textbf{Double DQN} --- separate action-selection and evaluation
        networks to reduce overestimation.
  \item \textbf{Dueling DQN} --- separate value $V(s)$ and advantage
        $A(s,a)$ streams.
  \item \textbf{PER} --- prioritised experience replay weighted by TD
        error.
\end{enumerate}

\textbf{Lagrangian constraint.}
A constrained variant adds $\lambda$ as a trainable multiplier updated
every 20 episodes via dual ascent:
$\lambda \leftarrow \operatorname{clip}\!\bigl(\lambda +
\eta(\bar u - B),\;0,\;\lambda_{\max}\bigr)$,
where $B$ is the allowed unserved budget and $\eta$ is the step size.

% =============================================================
\section{Baselines}

\textbf{Idle} --- battery never used.
\textbf{Rule} --- sells at peak, stores spare solar, keeps a 30\,\%
reserve, sheds during blackouts.
\textbf{Oracle} --- per-day perfect-foresight dynamic programming;
an upper-bound benchmark, not our agent.

% =============================================================
\section{Evaluation}
\label{sec:eval}

\textbf{Protocol.}
All agents are evaluated on test year 2023 using 200 fixed (day,
outage) scenarios drawn with seed 11.  Each DQN variant is trained
with 3 seeds (0--2); we report mean\,$\pm$\,std.  Unconstrained
variants use their final checkpoint; constrained models are shown
with both the final and the best-validation checkpoint (selected on
2022).

Return $G$ and the reported \emph{Return} column are both computed
with $\lambda{=}1$ for \emph{every} policy (constrained agents were
trained with larger $\lambda$, so they are measured on a common
yardstick rather than their own objective).
``Cost'' is energy purchases plus shedding, forced-cut and wear
penalties, and \emph{excludes} the critical-load penalty, which is
reported separately as unserved energy (kWh/day).

\textit{Note on randomness:} the (day, outage) scenarios are identical
across all policies; forecast and outage-warning noise is drawn from
each environment's own random stream and is therefore not bit-for-bit
identical across policies---but the scenario set is the same.

\begin{table}[h]
\centering
\caption{Test-year 2023 results (200 scenarios, 3 seeds,
  mean\,$\pm$\,std).
  ``Gap'' $= (\text{agent}-\text{rule})\,/\,(\text{oracle}-\text{rule})$.
  Cost excludes critical-load penalty (see text).}
\label{tab:results}
\setlength\tabcolsep{3.5pt}
\small
\begin{tabular}{lrrrr}
\toprule
Method & Return & Cost & Unserved & Gap \\
       &        & (Rs/d) & (kWh/d) & closed \\
\midrule
rule              & $-14.80$ & 14744 & 0.30 & 0\%          \\
oracle            & $-11.77$ & 11772 & 0.00 & 100\%        \\
\midrule
DQN               & $-12.44{\pm}0.07$ & 12319 & 0.60 & $78.0\%{\pm}2.4$ \\
+~Double          & $-12.36{\pm}0.09$ & 12272 & 0.44 & $80.6\%{\pm}2.9$ \\
+~Dueling         & $-12.41{\pm}0.13$ & 12308 & 0.50 & $79.0\%{\pm}4.3$ \\
+~PER             & $-12.65{\pm}0.10$ & 12486 & 0.82 & $71.1\%{\pm}3.3$ \\
Full (D2Q+PER)    & $-12.39{\pm}0.09$ & 12240 & 0.74 & $79.7\%{\pm}2.9$ \\
\midrule
constr.\ (final)    & $-13.02{\pm}0.39$ & 12683 & 1.67 & $59.0\%{\pm}12.9$ \\
constr.\ (best-val) & $-12.82{\pm}0.23$ & 12721 & 0.48 & $65.5\%{\pm}7.6$ \\
constr.\_v2 (final) & $-12.70{\pm}0.13$ & 12600 & 0.50 & $69.4\%{\pm}4.4$ \\
\bottomrule
\end{tabular}
\end{table}

\subsection*{Discussion}

All unconstrained agents beat the rule baseline (16--17\,\% lower
operating cost, 71--81\,\% of the oracle gap closed) but leave
slightly more critical load unserved than the rule.  Differences
among DQN variants are largely within seed noise; PER alone performs
slightly worse, likely due to instability early in training.

For the first constrained run, $\lambda$ reached its cap of 20 in all
three seeds and training was highly unstable (gap $59.0\%\pm12.9$).
For \texttt{constrained\_v2} ($\lambda_{\max}{=}5$, budget
$=0.5$\,kWh/day, $\eta{=}0.05$), $\lambda$ \emph{also} reached its cap
of 5 in all seeds; the tighter cap kept training more stable (gap
$69.4\%\pm4.4$) but unserved critical energy (0.50\,kWh/day) barely
met the budget, and cost performance remains below the unconstrained
full agent (79.7\%), suggesting the Lagrangian scheme trades cost
performance for only marginal blackout protection at this budget level.

% =============================================================
\section{Limitations}

All environment parameters (load, tariff, outages) are assumptions;
data come from one site with coarse satellite-derived solar irradiance;
three seeds is insufficient to detect small differences; and neither
constrained variant converged to a satisfactory unserved-energy budget.

% =============================================================
\section{Reproducibility}
{\small
\begin{verbatim}
pip install -r requirements.txt
python src/test_env.py   # sanity check
python src/demo.py       # demo agent
\end{verbatim}
}
\noindent
Approximate CPU runtime: $\sim$16\,min per 3\,000-episode run
(measured: 1.3\,min for 250 episodes, scaled $\times 12$).
All trained weights (\texttt{.pt}), logs (\texttt{.json}), and figures
(\texttt{.png}) are in \texttt{results/}.

\end{document}
